\section{AMC 12 12B 2026}
        \begin{enumerate}
        	\item (\textit{AMC 12 12B 2026}) A souvenir shop sells solid metal sculptures of Achilles in two sizes. One is $6$ cm tall and one is $8$ cm tall. Proportions between the two sculptures are preserved. The smaller sculpture has a mass of $486$ grams. What is the mass, in grams, of the larger sculpture?


 $\textbf{(A)}\   648  \qquad                 \textbf{(B)}\   864  \qquad                 \textbf{(C)}\   972  \qquad                 \textbf{(D)}\   1152  \qquad                 \textbf{(E)}\   1458$


	\item (\textit{AMC 12 12B 2026}) Bees occupy the $8$ corners of a $1 \times 2 \times 3$ rectangular box. The bees can fly inside the box, but cannot fly farther than $\frac{1}{2}$ unit from their respective corners. What is the volume, in cubic units, of the region inside the box that cannot be reached by any bee?


 $\textbf{(A)}\   6 - \frac{\pi}{3}  \qquad                 \textbf{(B)}\   6 - \frac{\pi}{4}  \qquad                 \textbf{(C)}\   6 - \frac{\pi}{6}  \qquad                 \textbf{(D)}\   6 - \frac{\pi}{8}  \qquad                 \textbf{(E)}\   6 - \frac{\pi}{12}$


	\item (\textit{AMC 12 12B 2026}) There are $2026$ equally spaced points that all lie on a circle. A line segment is drawn between two points if there is exactly one point between them along the circle. Into how many regions will this divide the interior of the circle? The figure below shows an example with $10$ equally spaced points, in which case there are $21$ interior regions.


 
                    \begin{figure}[H]
                        \centering
                        \includegraphics[width=0.5\textwidth]{Problem/AMC_12_12B_2026_3_Problem_Figure_1.png}
                    \end{figure}
                    
$\textbf{(A)}\   4052  \qquad                 \textbf{(B)}\   4053  \qquad                 \textbf{(C)}\   4054  \qquad                 \textbf{(D)}\   6079  \qquad                 \textbf{(E)}\   6081$


	\item (\textit{AMC 12 12B 2026}) For how many integers $n$ is

 \[\frac{n^{2}-1}{n+2}\]
also an integer?


 $\textbf{(A)}\   2  \qquad                 \textbf{(B)}\   3  \qquad                 \textbf{(C)}\   4  \qquad                 \textbf{(D)}\   6  \qquad                 \textbf{(E)}\   \text{infinitely many}$


	\item (\textit{AMC 12 12B 2026}) Let $a = \log_{2}27$, let $b = \log_{3}25$, and let $c = \log_{32}5$. What is $\frac{ab}{c}$?


 $\textbf{(A)}\   25  \qquad                 \textbf{(B)}\   27  \qquad                 \textbf{(C)}\   30  \qquad                 \textbf{(D)}\   32  \qquad                 \textbf{(E)}\   36$


	\item (\textit{AMC 12 12B 2026}) The integer $n$ satisfies the inequalities

 \[\frac{n}{100} < \frac{1}{5 \cdot 10} + \frac{1}{10 \cdot 15} + \frac{1}{15 \cdot 20} + \cdots + \frac{1}{90\cdot 95} + \frac{1}{95\cdot 100} < \frac{n+1}{100}.\]
What is $n$?


 $\textbf{(A)}\   2  \qquad                 \textbf{(B)}\   3  \qquad                 \textbf{(C)}\   4  \qquad                 \textbf{(D)}\   5  \qquad                 \textbf{(E)}\   6$


	\item (\textit{AMC 12 12B 2026}) Sharon tosses a fair coin six times and keeps track of how many heads and tails she has gotten. Her first toss is heads. What is the probability that after every toss Sharon's count of heads differs from her count of tails?


 $\textbf{(A)}\   \frac{1}{8}  \qquad                 \textbf{(B)}\   \frac{3}{16}  \qquad                 \textbf{(C)}\   \frac{1}{4}  \qquad                 \textbf{(D)}\   \frac{5}{16}  \qquad                 \textbf{(E)}\   \frac{3}{8}$


	\item (\textit{AMC 12 12B 2026}) A regular octahedron has $8$ equilateral triangle faces with edges of length $\sqrt{2}$. A plane is perpendicular to the line segment joining two opposite vertices of the octahedron and intersects that line segment $\frac{1}{2}$ unit from one of its endpoints. What is the area of the intersection of that plane and the octahedron?


 $\textbf{(A)}\   \frac{1}{4}  \qquad                 \textbf{(B)}\   \frac{1}{2}  \qquad                 \textbf{(C)}\   \frac{1}{2}\sqrt{2}  \qquad                 \textbf{(D)}\   \frac{3}{4}  \qquad                 \textbf{(E)}\   \sqrt{2}$


	\item (\textit{AMC 12 12B 2026}) Suppose that $a$ and $b$ are positive integers satisfying $20a + 26b = 2026$. What is the least possible value of $a+b$?


 $\textbf{(A)}\   71  \qquad                 \textbf{(B)}\   80  \qquad                 \textbf{(C)}\   83  \qquad                 \textbf{(D)}\   86  \qquad                 \textbf{(E)}\   88$


	\item (\textit{AMC 12 12B 2026}) Point $C$ lies on a circle with diameter $\overline{AB}$. Point $D$ lies on $\overline{AB}$, $14$ units from $A$, $24$ units from $B$, and $16$ units from $C$. What is the area, in square units, of $\triangle ABC$?


 $\textbf{(A)}\   186\sqrt{2}  \qquad                 \textbf{(B)}\   152\sqrt{3}  \qquad                 \textbf{(C)}\   264  \qquad                 \textbf{(D)}\   187\sqrt{2}  \qquad                 \textbf{(E)}\   108\sqrt{6}$


	\item (\textit{AMC 12 12B 2026}) A list contains $2026$ statements, each of which is either true or false. The first statement is "The number of false statements in this list is at least $1001$." The second statement is "The number of false statements in this list is at least $1002$." In general, the $k$th statement in the list is "The number of false statements in this list is at least $k+1000$." How many statements in the list are true?


 $\textbf{(A)}\   342  \qquad                 \textbf{(B)}\   513  \qquad                 \textbf{(C)}\   1013  \qquad                 \textbf{(D)}\   1026  \qquad                 \textbf{(E)}\   1513$


	\item (\textit{AMC 12 12B 2026}) Polly and Quentin are dancing. They start $1$ meter apart. Polly rotates $120^{\circ}$ clockwise around Quentin while he remains still, staying $1$ meter away throughout. Then Quentin rotates $120^{\circ}$ clockwise around Polly while she remains still, staying $1$ meter away throughout. They repeatedly do this pair of moves from their new positions. How many meters from where she started is Polly after they have each moved $10$ times?


 $\textbf{(A)}\   0  \qquad                 \textbf{(B)}\   \frac{\sqrt{3}}{2}  \qquad                 \textbf{(C)}\   1  \qquad                 \textbf{(D)}\   \sqrt{3}  \qquad                 \textbf{(E)}\   2$


	\item (\textit{AMC 12 12B 2026}) A jar contains $12$ green marbles and $8$ red marbles. In a game, Anika and Ben take turns drawing a random marble from the jar, with replacement. If Ben ever draws a red marble, he wins, and if Anika ever draws a green marble, she wins. Otherwise, the game continues with the next player's turn. The probability that Ben eventually wins is $A$ if Anika goes first and $B$ if Ben goes first. What is $B - A$?


 $\textbf{(A)}\   \frac{2}{25}  \qquad                 \textbf{(B)}\   \frac{4}{25}  \qquad                 \textbf{(C)}\   \frac{4}{19}  \qquad                 \textbf{(D)}\   \frac{6}{19}  \qquad                 \textbf{(E)}\   \frac{9}{20}$


	\item (\textit{AMC 12 12B 2026}) A line segment is drawn between two points on the graph of $y = x^{2}$ whose $x$-coordinates have been chosen independently and uniformly at random in the interval $[0,2]$. The probability that the $y$-coordinate of the midpoint of the segment is greater than $2$ lies in which of the following intervals?


 $\textbf{(A)}\   (0.20, 0.21)  \qquad                 \textbf{(B)}\   (0.21, 0.22)  \qquad                 \textbf{(C)}\   (0.22, 0.23)  \qquad                 \textbf{(D)}\   (0.23, 0.24)  \qquad                 \textbf{(E)}\   (0.24, 0.25)$


	\item (\textit{AMC 12 12B 2026}) What is the value of

 \[\log_{2}\left( \sin\left( \frac{\pi}{8} \right)  \right) + \log_{2}\left( \sin\left( \frac{3\pi}{8} \right)  \right) + \log_{2}\left( \sin\left( \frac{5\pi}{8} \right)  \right) + \log_{2}\left( \sin\left( \frac{7\pi}{8} \right)  \right)?\]
$\textbf{(A)}\   -3  \qquad                 \textbf{(B)}\   -2\sqrt{2}  \qquad                 \textbf{(C)}\   -2  \qquad                 \textbf{(D)}\   -\frac{3}{2}  \qquad                 \textbf{(E)}\   -\sqrt{2}$


	\item (\textit{AMC 12 12B 2026}) Annie, Bobby, and Casey each receive an integer, not necessarily distinct, from $1$ through $10$. In how many ways can the numbers be given so that their sum is divisible by $7$?


 $\textbf{(A)}\   139  \qquad                  \textbf{(B)}\   140  \qquad                 \textbf{(C)}\   142  \qquad                 \textbf{(D)}\   143  \qquad                 \textbf{(E)}\   145$


	\item (\textit{AMC 12 12B 2026}) Belinda follows a path on the plane that is entirely contained within the square region $0 \leq x, y \leq 4$. Call her path \textit{dignified} if and only if all of the following conditions are satisfied.


 
            \begin{itemize}
        \item Each step is one unit up, down, or to the right.

        \item The path does not revisit any vertex.

        \item The path starts at $(a, 0)$ and ends at $(b,0)$, where a and $b$ are integers with $a < b$, and touches the $x$-axis only at these two points.
            \end{itemize}
            The following figure shows some examples of dignified and not dignified paths.


 
                    \begin{figure}[H]
                        \centering
                        \includegraphics[width=0.5\textwidth]{Problem/AMC_12_12B_2026_17_Problem_Figure_1.png}
                    \end{figure}
                    


 How many dignified paths can Belinda follow?


 $\textbf{(A)}\   174  \qquad                 \textbf{(B)}\   340  \qquad                 \textbf{(C)}\   448  \qquad                 \textbf{(D)}\   512  \qquad                 \textbf{(E)}\   970$


	\item (\textit{AMC 12 12B 2026}) A ribbon wraps around a pole in such a way that equal amounts of area are covered and not covered by the ribbon. The width of the ribbon is $4$ cm, and the circumference of the pole is $10$ cm. The figure below shows a possible front view of a portion of the pole.


 
                    \begin{figure}[H]
                        \centering
                        \includegraphics[width=0.5\textwidth]{Problem/AMC_12_12B_2026_18_Problem_Figure_1.png}
                    \end{figure}
                    


 In centimeters, how much of the pole height is needed for the ribbon to wrap around exactly once?


 $\textbf{(A)}\   \frac{15}{2}  \qquad                 \textbf{(B)}\   8  \qquad                 \textbf{(C)}\   10  \qquad                 \textbf{(D)}\   \frac{25}{2}  \qquad                 \textbf{(E)}\   \frac{40}{3}$


	\item (\textit{AMC 12 12B 2026}) Starting with a list of three numbers, form a sequence in which each subsequent term is the average of the previous three terms. For example, starting with the sequence $(2,4,9)$, the next two terms are $\frac{2+4+9}{3} = 5$ and $\frac{4+9+5}{3} = 6$. Which of the following initial three terms will produce a sequence that has the greatest 2026th term?


 $\textbf{(A)}\   (1,2,3)  \qquad                 \textbf{(B)}\   (2,1,1)  \qquad                 \textbf{(C)}\   (3,4,8)  \qquad                 \textbf{(D)}\   (4,4,7)  \qquad                 \textbf{(E)}\   (4,5,6)$


	\item (\textit{AMC 12 12B 2026}) Which of the following statements most accurately describes the graph of the relationship

 \[\sin\left( x^{2} \right) = \cos\left( y^{2} \right)\]
$\textbf{(A)}\   \text{the set of all points } \left( \sqrt{ \frac{m}{\pi}}, \sqrt{\frac{n + \frac{1}{2}}{\pi}} \right), \text{ where } m \text{ and } n \text{ are nonnegative integers}$
$\textbf{(B)}\   \text{the set of all lines } y = \pm x + \frac{\pi}{2}, \text{ where } k \text{ is an integer}$
$\textbf{(C)}\   \text{the set of all circles } x^{2} + y^{2} = \frac{\pi}{2} + 2 \pi n, \text{ where n is a nonnegative integer}$
$\textbf{(D)}\   \text{a set of infinitely many concentric circles and infinitely many hyperbolas}$
$\textbf{(E)}\   \text{a set of finitely many ellipses centered at the origin and infinitely many parabolas}$


	\item (\textit{AMC 12 12B 2026}) Let $T_{n}$ be the $n$th triangular number, that is, $T_{1} = 1$, $T_{2} = 1+2 = 3$, $T_{3} = 1+2+3 = 6$, and so on. The value of

 \[\sum_{k=1}^{\infty} T_{k} \cdot \left( \frac{1}{3} \right)^{k}\]
can be expressed as $\frac{p}{q}$, where $p$ and $q$ are relatively prime positive integers. What is $p+q$?


 $\textbf{(A)}\   14  \qquad                 \textbf{(B)}\   15  \qquad                 \textbf{(C)}\   16  \qquad                 \textbf{(D)}\   17  \qquad                 \textbf{(E)}\   18$


	\item (\textit{AMC 12 12B 2026}) Rectangular card $ABCD$ has $AB = 1$ and $AC = 2$. As shown in the figure, a second card, $AEFG$, congruent to the first, is stacked onto it so that $\overline{AG}$ lies on $\overline{AC}$ and the two cards share corner $A$.


 
                    \begin{figure}[H]
                        \centering
                        \includegraphics[width=0.5\textwidth]{Problem/AMC_12_12B_2026_22_Problem_Figure_1.png}
                    \end{figure}
                    


 Additional cards are added, following the same pattern (with shared corner $A$), creating a fan of cards. The process stops when a new card fits exactly atop $ABCD$. The shape then created is a nonconvex polygon with rotational symmetry. What is the area of this polygon?


 $\textbf{(A)}\   20\sqrt{3} - 24  \qquad                 \textbf{(B)}\   34\sqrt{3} - 48  \qquad                 \textbf{(C)}\   6 + 3\sqrt{3}  \qquad                 \textbf{(D)}\   12  \qquad                 \textbf{(E)}\   8\sqrt{3}$


	\item (\textit{AMC 12 12B 2026}) For positive integers $a$ and $b$, let $R(a, b)$ denote the remainder when $a$ is divided by $b$. Suppose $N$ is the least positive integer such that

 \[R(N,2) \ge 1, \quad R(N,4) \ge 2, \quad R(N,8) \ge 4, \quad \dots, \quad R\left(N,2^{200}\right) \ge 2^{199}.\]
What are the last three digits of $N$?


 $\textbf{(A)}\   125  \qquad                 \textbf{(B)}\   249  \qquad                 \textbf{(C)}\   375  \qquad                 \textbf{(D)}\   511  \qquad                 \textbf{(E)}\   875$


	\item (\textit{AMC 12 12B 2026}) Let $\triangle ABC$ be a scalene triangle with incenter $I$. (Recall that the incenter is the intersection of the three angle bisectors.) The distances from $I$ to the midpoints of the sides of $\triangle ABC$ are $\sqrt{11}$, $\sqrt{14}$, and $\sqrt{19}$. What is the sum of the possible areas of the circle inscribed in $\triangle ABC$?


 $\textbf{(A)}\   10 \pi  \qquad                 \textbf{(B)}\   \frac{44 \pi}{3}  \qquad                 \textbf{(C)}\   \frac{58 \pi}{3}  \qquad                 \textbf{(D)}\   \frac{88 \pi}{3}  \qquad                 \textbf{(E)}\   30 \pi$


	\item (\textit{AMC 12 12B 2026}) A \textit{finite integer tower} is an expression of the form

 \[a_{1}^{a_{2}^{a_{3}^{\iddots^{a_{n}}}}}\]
where $n \geq 1$ and $a_{1}$, $a_2$, $\dots$, $a_{n}$ are positive integers greater than $1$. (Recall that exponentiation is performed from the top down. For example, $3^{2^{3}}$ equals $3^{8}$, not $9^{3}$.) Two towers are considered different if they have a different number of layers (values of $n$) or at least one pair of numbers in corresponding positions differ.


 Define $f(x)$ to be the number of different finite integer towers whose value is $x$. (For example, $f(16)=4$ because $4$ towers evaluate to $16$, namely $16$, $2^{4}$, $4^2$, and $2^{2^{2}}$.) What is $f(2^{8192})$?


 $\textbf{(A)}\   17  \qquad                 \textbf{(B)}\   32  \qquad                 \textbf{(C)}\   42  \qquad                 \textbf{(D)}\   43  \qquad                 \textbf{(E)}\   47$


\end{enumerate}